\chapter{Global Error Handling}

Scattered \texttt{try/catch} blocks that each format their own error
response lead to inconsistent API output -- one endpoint returns
\texttt{\{ error: "..." \}}, another returns \texttt{\{ message: "..." \}},
and a third crashes the process entirely on an unhandled rejection. A single
global error handling middleware fixes all three problems at once: every
error, from anywhere in the app, funnels through one place that shapes the
response consistently.

\begin{diagrambox}[Error Handling Flow]
\begin{verbatim}
 Controller
    |  something goes wrong
    v
 throw new AppError("Task not found", 404)
    |
    v
 catch block -> next(error)
    |
    v
 (skips all remaining normal middleware/routes)
    |
    v
 Global Error Middleware  (app.use((err, req, res, next) => ...))
    |
    v
 JSON Response  { success: false, message }
\end{verbatim}
\end{diagrambox}

\section{A Custom AppError Class}

Plain \texttt{Error} objects don't carry an HTTP status code. A small
subclass fixes that, letting any part of the app throw an error that already
knows how it should be reported to the client.

\begin{lstlisting}[language=JavaScript]
// utils/AppError.js
class AppError extends Error {
  constructor(message, statusCode) {
    super(message);
    this.statusCode = statusCode;
    this.isOperational = true; // distinguishes expected errors from bugs
  }
}

export default AppError;
\end{lstlisting}

\begin{lstlisting}[language=JavaScript]
// controllers/taskController.js
import AppError from "../utils/AppError.js";

export const getTaskById = async (req, res, next) => {
  try {
    const task = await Task.findById(req.params.id);
    if (!task) {
      throw new AppError("Task not found", 404);
    }
    res.status(200).json({ success: true, data: task });
  } catch (err) {
    next(err);
  }
};
\end{lstlisting}

\section{Avoiding Repetitive try/catch: asyncHandler}

Writing \texttt{try \{ ... \} catch (err) \{ next(err) \}} in every single
controller is repetitive and easy to forget. A small wrapper function
eliminates it by catching any rejected promise automatically.

\begin{lstlisting}[language=JavaScript]
// utils/asyncHandler.js
const asyncHandler = (fn) => (req, res, next) => {
  Promise.resolve(fn(req, res, next)).catch(next);
};

export default asyncHandler;
\end{lstlisting}

\begin{lstlisting}[language=JavaScript]
// controllers/taskController.js
import asyncHandler from "../utils/asyncHandler.js";
import AppError from "../utils/AppError.js";

export const getTaskById = asyncHandler(async (req, res) => {
  const task = await Task.findById(req.params.id);
  if (!task) {
    throw new AppError("Task not found", 404);
  }
  res.status(200).json({ success: true, data: task });
});
\end{lstlisting}

Any thrown error, or any rejected promise from an \texttt{await} call inside
the wrapped function, is automatically passed to \texttt{next(err)} -- the
controller body no longer needs its own \texttt{try/catch} at all.

\section{The Global Error Middleware}

Express recognizes a middleware with four parameters
\texttt{(err, req, res, next)} as an error handler. It must be registered
\emph{after} all routes, so anything that calls \texttt{next(err)} anywhere
in the app eventually reaches it.

\begin{lstlisting}[language=JavaScript]
// middleware/errorHandler.js
export const errorHandler = (err, req, res, next) => {
  let statusCode = err.statusCode || 500;
  let message = err.message || "Internal Server Error";

  // Mongoose: invalid ObjectId format, e.g. Task.findById("not-an-id")
  if (err.name === "CastError") {
    statusCode = 400;
    message = `Invalid ${err.path}: ${err.value}`;
  }

  // Mongoose: schema validation failed (required field, enum, etc.)
  if (err.name === "ValidationError") {
    statusCode = 400;
    message = Object.values(err.errors).map((e) => e.message).join(", ");
  }

  // MongoDB: duplicate key on a unique field, e.g. email already exists
  if (err.code === 11000) {
    statusCode = 409;
    const field = Object.keys(err.keyValue)[0];
    message = `${field} already exists`;
  }

  res.status(statusCode).json({
    success: false,
    message,
  });
};
\end{lstlisting}

\begin{lstlisting}[language=JavaScript]
// app.js
import { errorHandler } from "./middleware/errorHandler.js";

app.use("/api/auth", authRoutes);
app.use("/api/tasks", taskRoutes);

// must be registered after all routes
app.use(errorHandler);
\end{lstlisting}

\begin{warnbox}[WARNING]
Middleware order matters. If \texttt{errorHandler} is registered before the
routes, Express never reaches it when a route calls \texttt{next(err)} --
the request hangs or Express falls back to its own default (and much less
helpful) error page.
\end{warnbox}

\section{Mapping Mongoose-Specific Errors}

\begin{center}
\begin{tabularx}{\textwidth}{lXl}
\toprule
\textbf{Error} & \textbf{Cause} & \textbf{Mapped status} \\
\midrule
\texttt{CastError} & Malformed value for a typed field, most often an invalid \texttt{ObjectId} in a URL param. & 400 \\
\texttt{ValidationError} & A document failed schema rules (\texttt{required}, \texttt{enum}, \texttt{minlength}, etc.) on save/create. & 400 \\
Code \texttt{11000} & Duplicate key write against a field with a \texttt{unique} index (e.g. registering an email that already exists). & 409 \\
\bottomrule
\end{tabularx}
\end{center}

\begin{tipbox}[TIP]
Because these checks live inside the single \texttt{errorHandler}, no
individual controller needs to know about Mongoose error internals --
\texttt{Task.findById(badId)} throwing a \texttt{CastError} is turned into a
clean \texttt{400 \{ message: "Invalid \_id: ..." \}} response automatically,
regardless of which controller triggered it.
\end{tipbox}

\whatwhyhow
{A single error-handling middleware, registered last in app.js, that every thrown or forwarded error eventually reaches and formats into a consistent JSON response.}
{Without it, error handling is duplicated and inconsistent across every controller, and raw Mongoose/MongoDB errors leak internal details or default to an unhelpful 500 instead of a meaningful status code.}
{Throw an \texttt{AppError(message, statusCode)} (or let asyncHandler catch any rejected promise), forward it with \texttt{next(err)}, and let \texttt{app.use((err, req, res, next) => ...)} map it -- including special cases like \texttt{CastError}, \texttt{ValidationError}, and duplicate key code 11000 -- to the right status and message.}
